\documentclass[letterpaper,10pt]{article}

\usepackage[letterpaper,margin=0.55in]{geometry}
\usepackage[english]{babel}
\usepackage{enumitem}
\usepackage[hidelinks]{hyperref}
\usepackage{titlesec}
\usepackage{tabularx}
\usepackage{xcolor}
\input{glyphtounicode}
\pdfgentounicode=1

\pagestyle{empty}
\setlength{\parindent}{0pt}
\setlength{\tabcolsep}{0pt}
\raggedright

\titleformat{\section}{\vspace{-2pt}\bfseries\large}{}{0pt}{}[\vspace{1pt}\titlerule]
\titlespacing*{\section}{0pt}{7pt}{5pt}

\newcommand{\entry}[4]{%
  \begin{tabular*}{\textwidth}{l@{\extracolsep{\fill}}r}
    \textbf{#1} & #2 \\
    \textit{#3} & \textit{#4}
  \end{tabular*}\vspace{-4pt}
}
\newcommand{\bullets}{\begin{itemize}[leftmargin=0.18in,topsep=3pt,itemsep=2pt,parsep=0pt,partopsep=0pt]}
\newcommand{\finishbullets}{\end{itemize}\vspace{-3pt}}
\newcommand{\pub}[1]{\item #1}

\begin{document}

\begin{center}
  {\LARGE\textbf{Haoyu Dong}}\\[4pt]
  \href{mailto:haoyu.dong151@duke.edu}{haoyu.dong151@duke.edu} $\vert$
  +1 858 500-2825 $\vert$
  \href{https://haoyudong-97.github.io/}{Homepage} $\vert$
  \href{https://www.linkedin.com/in/haoyu-dong-63a801b9/}{LinkedIn} $\vert$
  \href{https://scholar.google.com/citations?hl=en\&user=eZVEUCIAAAAJ\&view_op=list_works}{Google Scholar}
\end{center}
\vspace{-8pt}

\section{Research Profile}
Ph.D. candidate in Electrical \& Computer Engineering at Duke University, specializing in foundation models, LLM agents, multimodal learning, and data-centric AI. Industry research experience at Meta and Siemens Healthineers, developing algorithm-evolution methods for LLM agents and in-context learning systems for 3D medical image segmentation. Author of 15+ peer-reviewed publications in medical AI, computer vision, and multimodal learning.

\section{Industry Experience}
\entry{Meta}{May 2026 -- Present}{Software Engineer Intern}{Menlo Park, CA}
\bullets
  \item Improved Meta-Harness's algorithm-evolution pipeline by making agent trajectories evolvable and dynamically adapting the training set to strengthen the optimization signal. Increased performance from \textbf{49\% to 62\%} on math-solving tasks (\textbf{+13 percentage points}) and from \textbf{68\% to 71.3\%} on Terminal-Bench 2 (\textbf{+3.3 percentage points}).
\finishbullets

\entry{Siemens Healthineers}{Jun. 2025 -- Aug. 2025}{AI Research Intern}{Princeton, NJ}
\bullets
  \item Developed a training-free in-context segmentation framework with fixed support-query prompting; built preprocessing and evaluation infrastructure spanning \textbf{14K MRI/CT volumes}. Co-inventor on the resulting \textbf{patent-pending} technology assigned to Siemens Medical Solutions USA, Inc.
  \item With only five supports per task and no retraining, achieved \textbf{52.0\% Dice} versus \textbf{60.0\%} for a supervised nnU-Net on an internal test set, and \textbf{54.9\%} versus \textbf{56.9\%} on the AMOS CT validation set.
\finishbullets

\entry{Gradient Health}{May 2020 -- Sep. 2020}{Fellowship Recipient}{Durham, NC}
\bullets
  \item Built a private multimodal dataset of diagnostic images and paired radiology reports from \textbf{2.6M patients}; developed a language-informed BYOL pretraining pipeline that improved accuracy by \textbf{3 percentage points} over a ResNet baseline.
\finishbullets

\section{Selected Research Experience}
\entry{Mazurowski's Lab, Duke University}{Jul. 2022 -- Present}{Graduate Research Associate}{Durham, NC}
\bullets
  \item Pretrained \textbf{MRI-CORE} on \textbf{7M MRI slices} using modern self-supervised objectives, improving few-shot downstream segmentation by \textbf{5 Dice points} over SAM; currently developing a 3D multimodal breast MRI foundation model that aligns volumetric imaging with radiology reports.
  \item Led a broad evaluation of promptable segmentation across \textbf{19 datasets, 5 modalities, and 18 body locations}, establishing data-efficient interactive workflows that transfer across scanners and sites (\textit{Medical Image Analysis}, 2023).
  \item Developed \textbf{InTEnt}, a single-image test-time adaptation method using integrated entropy weighting to adjust batch-normalization statistics, improving out-of-domain segmentation by \textbf{2.9 Dice points} over the strongest comparator (CVPR 2024 Workshop Oral).
  \item Co-developed foundation-model adaptation guidelines spanning prompt design, parameter-efficient adaptation, and light fine-tuning, and contributed to universal anatomical segmentation models across locations and acquisition protocols.
\finishbullets

\entry{Mazurowski's Lab, Duke University}{Jun. 2021 -- Jun. 2022}{Associate in Research}{Durham, NC}
\bullets
  \item Proposed \textbf{SWSSL} for anomaly detection in high-resolution breast tomosynthesis and chest X-ray, surpassing prior state of the art by \textbf{5.0 AUROC points} on tomosynthesis (\textit{IEEE Transactions on Medical Imaging}, 2023).
\finishbullets

\entry{Sapiro's Lab, Duke University}{Sep. 2019 -- May 2021}{Graduate Research Assistant}{Durham, NC}
\bullets
  \item Built a visual-semantic space for text-conditioned image retrieval using iterative point alignment and manifold-distance search, improving recall by \textbf{2 points} on FashionIQ, \textbf{5 points} on CSS, and \textbf{30 points} on a new dataset (CVPR 2021 Workshop).
\finishbullets

\section{Selected Publications}
\begin{itemize}[leftmargin=0.18in,topsep=2pt,itemsep=4pt,parsep=0pt]
  \pub{\textbf{Haoyu Dong}, Yuwen Chen, Hanxue Gu, Nicholas Konz, Yaqian Chen, Qihang Li, and Maciej A. Mazurowski. ``MRI-CORE: A Foundation Model for Magnetic Resonance Imaging.'' \textit{Preprint}, 2025.}
  \pub{\textbf{Haoyu Dong}, Han Liu, Riqiang Gao, Bogdan Georgescu, Guillaume Chabin, Jianing Wang, and Sasa Grbi\'c. ``Revisiting In-Context Learning for 3D CT Segmentation of Unseen Anatomies.'' \textit{Submitting}.}
\end{itemize}
\newpage
\section{Selected Publications (continued)}
\begin{itemize}[leftmargin=0.18in,topsep=2pt,itemsep=4pt,parsep=0pt]
  \pub{Hanxue Gu*, \textbf{Haoyu Dong*}, Pedro R. A. S. Bassi, Wenxuan Li, Alan Yuille, Zongwei Zhou, Kang Wang, and Yang Yang. ``Reports Are Not Captions: Understanding and Designing Report Supervision for Label-Efficient 3D Lesion Segmentation.'' \textit{Submitting}.}
  \pub{Hanxue Gu*, \textbf{Haoyu Dong*}, Jichen Yang, and Maciej A. Mazurowski. ``How to build the best medical image segmentation algorithm using foundation models: a comprehensive empirical study with Segment Anything Model.'' \textit{Machine Learning for Biomedical Imaging}, 2025.}
  \pub{\textbf{Haoyu Dong}, Nicholas Konz, Hanxue Gu, and Maciej A. Mazurowski. ``Medical image segmentation with InTEnt: Integrated entropy weighting for single image test-time adaptation.'' \textit{CVPR Workshops}, 2024.}
  \pub{Maciej A. Mazurowski, \textbf{Haoyu Dong}, Hanxue Gu, Jichen Yang, Nicholas Konz, and Yixin Zhang. ``Segment anything model for medical image analysis: an experimental study.'' \textit{Medical Image Analysis}, 2023.}
  \pub{\textbf{Haoyu Dong}, Yifan Zhang, Hanxue Gu, Nicholas Konz, Yixin Zhang, and Maciej A. Mazurowski. ``SWSSL: Sliding window-based self-supervised learning for anomaly detection in high-resolution images.'' \textit{IEEE Transactions on Medical Imaging}, 2023.}
\end{itemize}
\vspace{-5pt}
\small * Equal contribution. Full publication list available on \href{https://scholar.google.com/citations?hl=en\&user=eZVEUCIAAAAJ\&view_op=list_works}{Google Scholar}.

\section{Education}
\entry{Duke University}{Expected Jan. 2027}{Ph.D., Electrical \& Computer Engineering}{Durham, NC}
\vspace{6pt}

\entry{Duke University}{Jun. 2021}{M.S., Computer Science}{Durham, NC}
\vspace{8pt}

\entry{University of California, San Diego}{Jun. 2019}{B.S., Mathematics--Computer Science; B.S., Cognitive Science}{La Jolla, CA}

\section{Technical Skills}
\textbf{Programming:} Python, C++\\[2pt]
\textbf{Machine Learning:} PyTorch, distributed training, reinforcement learning, multimodal modeling\\[2pt]
\textbf{Infrastructure:} Linux, Git, Docker, Slurm, cloud platforms

\section{Academic Service}
\textbf{Conference Reviewer:} ICLR, NeurIPS, ICML, CVPR, ICCV, ECCV, MICCAI, MIDL\\[2pt]
\textbf{Journal Reviewer:} IEEE TMI (Distinguished Reviewer, Silver Level 2023--24), IEEE JBHI, IEEE TNNLS, IEEE TBME, Medical Image Analysis, npj Digital Medicine, npj Artificial Intelligence

\end{document}
